\subsection{Example: A Picture Language}
\label{Section 2.2.4}

This section presents a simple language for drawing pictures that illustrates
the power of data abstraction and closure, and also exploits higher-order
procedures in an essential way.  The language is designed to make it easy to
experiment with patterns such as the ones in \link{Figure 2.9}, which are
composed of repeated elements that are shifted and scaled.\footnote{The picture
language is based on the language Peter Henderson created to construct images
like M.C. Escher's ``Square Limit'' woodcut (see \link{Henderson 1982}).  The
woodcut incorporates a repeated scaled pattern, similar to the arrangements
drawn using the \code{square\_limit} procedure in this section.} In this
language, the data objects being combined are represented as procedures rather
than as list structure.  Just as \code{cons}, which satisfies the closure
property, allowed us to easily build arbitrarily complicated list structure,
the operations in this language, which also satisfy the closure property, allow
us to easily build arbitrarily complicated patterns.


\subsubsection*{The picture language}

When we began our study of programming in \link{Section 1.1}, we emphasized the
importance of describing a language by focusing on the language's primitives,
its means of combination, and its means of abstraction.  We'll follow that
framework here.

% figure 2.9
\begin{imageFigure}{A design generated with the picture language:
  \code{square\_limit(wave, 3)}, drawn by the code of this section.}
  \generatedFigure[76mm]{ch2/wave-square-limit}
\end{imageFigure}

% figure 2.10
\begin{imageFigure}
  {%
    The image produced by the \code{wave} painter with respect to the unit
    square.  The same painter draws a different image in a different frame.
  }%
  \generatedFigure[42mm]{ch2/wave}
\end{imageFigure}

Part of the elegance of this picture language is that there is only one kind of
element, called a \newterm{painter}.  A painter draws an image that is shifted
and scaled to fit within a designated parallelogram-shaped frame.  For example,
there's a primitive painter we'll call \code{wave} that makes a crude line
drawing, as shown in \link{Figure 2.10}.  The actual shape of the drawing
depends on the frame---the same \code{wave} painter draws a different image
with respect to a different frame.  Painters can be more elaborate than this:
the original has a primitive painter called \code{rogers} that paints a picture
of \acronym{MIT}'s founder, William Barton Rogers.\footnote{%
William Barton Rogers (1804-1882) was the founder
and first president of \acronym{MIT}. A geologist and talented teacher, he
taught at William and Mary College and at the University of Virginia.  In 1859
he moved to Boston, where he had more time for research, worked on a plan for
establishing a ``polytechnic institute,'' and served as Massachusetts's first
State Inspector of Gas Meters.

When \acronym{MIT} was established in 1861, Rogers was elected its first
president.  Rogers espoused an ideal of ``useful learning'' that was different
from the university education of the time, with its overemphasis on the
classics, which, as he wrote, ``stand in the way of the broader, higher and
more practical instruction and discipline of the natural and social sciences.''
This education was likewise to be different from narrow trade-school education.
In Rogers's words:

\begin{niceQuote}
The world-enforced distinction between the practical and the scientific worker
is utterly futile, and the whole experience of modern times has demonstrated
its utter worthlessness.
\end{niceQuote}

Rogers served as president of \acronym{MIT} until 1870, when he resigned due to
ill health.  In 1878 the second president of \acronym{MIT}, John Runkle,
resigned under the pressure of a financial crisis brought on by the Panic of
1873 and strain of fighting off attempts by Harvard to take over \acronym{MIT}.
Rogers returned to hold the office of president until 1881.

Rogers collapsed and died while addressing \acronym{MIT}'s graduating class at
the commencement exercises of 1882.  Runkle quoted Rogers's last words in a
memorial address delivered that same year:

\begin{niceQuote}
``As I stand here today and see what the Institute is, \( \dots \) I call to
mind the beginnings of science. I remember one hundred and fifty years ago
Stephen Hales published a pamphlet on the subject of illuminating gas, in which
he stated that his researches had demonstrated that 128 grains of bituminous
coal -- '' ``Bituminous coal,'' these were his last words on earth. Here he
bent forward, as if consulting some notes on the table before him, then slowly
regaining an erect position, threw up his hands, and was translated from the
scene of his earthly labors and triumphs to ``the tomorrow of death,'' where
the mysteries of life are solved, and the disembodied spirit finds unending
satisfaction in contemplating the new and still unfathomable mysteries of the
infinite future.
\end{niceQuote}

In the words of Francis A. Walker (\acronym{MIT}'s third president):

\begin{niceQuote}
All his life he had borne himself most faithfully and heroically, and he died
as so good a knight would surely have wished, in harness, at his post, and in
the very part and act of public duty.
\end{niceQuote}
}
\begin{buildFromScarcity}
Our painters draw in straight lines and nothing else, so \code{rogers} does not
come with us.  That is a loss worth naming, because the original is making a
point with it: a painter that paints a photograph and a painter that draws a
stick figure are the same kind of thing, and everything in this section works
on both without ever knowing which it has.  The argument survives---it is about
the interface and not the images---but the reader should know that the picture
language is capable of more than the pictures we shall draw with it.

The reason for the loss is the same one that runs through this section.  The
original assumes a graphics system and never says what it is, because in 1985
that depended on the machine in front of you.  We cannot leave it unsaid: our
\code{wave} has to be something the reader can run, something that can be
printed on this page, and something that can be drawn on a screen.  A sequence
of line segments satisfies all three at once.  What our painters produce, and
how it reaches the page, is the subject of the next few pages.
\end{buildFromScarcity}

To combine images, we use various operations that construct new painters from
given painters.  For example, the \code{beside} operation takes two painters
and produces a new, compound painter that draws the first painter's image in
the left half of the frame and the second painter's image in the right half of
the frame.  Similarly, \code{below} takes two painters and produces a compound
painter that draws the first painter's image below the second painter's image.
Some operations transform a single painter to produce a new painter.  For
example, \code{flip\_vert} takes a painter and produces a painter that draws
its image upside-down, and \code{flip\_horiz} produces a painter that draws the
original painter's image left-to-right reversed.

\link{Figure 2.11} shows the drawing of a painter called
\code{wave4} that is built up in two stages starting from \code{wave}:

\begin{codeBlock}
>>> wave2 = beside(wave, flip_vert(wave))
>>> wave4 = below(wave2, wave2)
\end{codeBlock}

% figure 2.12
\begin{imageFigure}
  {%
    Creating a complex figure, starting from the \code{wave} painter of
    \link{Figure 2.10}: \code{wave2} on the left, \code{wave4} on the right.
  }%
  \generatedFigure[38mm]{ch2/wave-beside}\hspace{6mm}%
  \generatedFigure[38mm]{ch2/wave-flipped-pairs}
\end{imageFigure}

In building up a complex image in this manner we are exploiting the fact that
painters are closed under the language's means of combination.  The
\code{beside} or \code{below} of two painters is itself a painter; therefore,
we can use it as an element in making more complex painters.  As with building
up list structure using \code{cons}, the closure of our data under the means of
combination is crucial to the ability to create complex structures while using
only a few operations.

Once we can combine painters, we would like to be able to abstract typical
patterns of combining painters.  We will implement the painter operations as
Python procedures.  This means that we don't need a special abstraction
mechanism in the picture language: since the means of combination are ordinary
Python procedures, we automatically have the capability to do anything with
painter operations that we can do with procedures.  For example, we can
abstract the pattern in \code{wave4} as

\begin{codeBlock}
>>> def flipped_pairs(painter):
...     painter2 = beside(painter, flip_vert(painter))
...     return below(painter2, painter2)
\end{codeBlock}

\noindent
and define \code{wave4} as an instance of this pattern:

\begin{codeBlock}
>>> wave4 = flipped_pairs(wave)
\end{codeBlock}

% figure 2.13
\begin{imageFigure}{Recursive plans for \code{right\_split} and
  \code{corner\_split}.}
  \includegraphics[width=111mm]{assets/figures/chapter_2/figure_2.13a.pdf}
\end{imageFigure}

\noindent
We can also define recursive operations.  Here's one that makes painters split
and branch towards the right as shown in \link{Figure 2.12}
and \link{Figure 2.13}:

\begin{codeBlock}
>>> def right_split(painter, n):
...     if n == 0:
...         return painter
...     else:
...         smaller = right_split(painter, n - 1)
...         return beside(painter, below(smaller, smaller))
\end{codeBlock}

\noindent
We can produce balanced patterns by branching upwards as well as towards the
right (see exercise \link{Exercise 2.44} and figures \link{Figure 2.12} and
\link{Figure 2.13}):

\begin{codeBlock}
>>> def corner_split(painter, n):
...     if n == 0:
...         return painter
...     else:
...         up = up_split(painter, n - 1)
...         right = right_split(painter, n - 1)
...         top_left = beside(up, up)
...         bottom_right = below(right, right)
...         corner = corner_split(painter, n - 1)
...         return beside(below(painter, top_left),
...                       below(bottom_right, corner))
\end{codeBlock}

\noindent
By placing four copies of a \code{corner\_split} appropriately, we obtain a
pattern called \code{square\_limit}, whose application to \code{wave} and
\code{rogers} is shown in \link{Figure 2.9}:

\begin{codeBlock}
>>> def square_limit(painter, n):
...     quarter = corner_split(painter, n)
...     half = beside(flip_horiz(quarter), quarter)
...     return below(flip_vert(half), half)
\end{codeBlock}

\begin{exerciseGroup}
\phantomsection\label{Exercise 2.44}
\exerciseHeading{Exercise 2.44:} Define the procedure
\code{up\_split} used by \code{corner\_split}.  It is similar to
\code{right\_split}, except that it switches the roles of \code{below} and
\code{beside}.
\end{exerciseGroup}

% figure 2.14
\begin{imageFigure}
  {%
    The recursive operations \code{right\_split} (left) and
    \code{corner\_split} (right) applied to \code{wave}, both at \code{n = 3}.
    Combining four \code{corner\_split} figures produces the symmetric
    \code{square\_limit} design of \link{Figure 2.9}.
  }%
  \generatedFigure[44mm]{ch2/wave-right-split}\hspace{5mm}%
  \generatedFigure[50mm]{ch2/wave-corner-split}
\end{imageFigure}

\subsubsection*{Higher-order operations}

In addition to abstracting patterns of combining painters, we can work at a
higher level, abstracting patterns of combining painter operations.  That is,
we can view the painter operations as elements to manipulate and can write
means of combination for these elements---procedures that take painter
operations as arguments and create new painter operations.

For example, \code{flipped\_pairs} and \code{square\_limit} each arrange four
copies of a painter's image in a square pattern; they differ only in how they
orient the copies.  One way to abstract this pattern of painter combination is
with the following procedure, which takes four one-argument painter operations
and produces a painter operation that transforms a given painter with those
four operations and arranges the results in a square.  \code{tl}, \code{tr},
\code{bl}, and \code{br} are the transformations to apply to the top left copy,
the top right copy, the bottom left copy, and the bottom right copy,
respectively.

\begin{codeBlock}
>>> def square_of_four(tl, tr, bl, br):
...     def combined(painter):
...         top = beside(tl(painter), tr(painter))
...         bottom = beside(bl(painter), br(painter))
...         return below(bottom, top)
...     return combined
\end{codeBlock}

\noindent
Then \code{flipped\_pairs} can be defined in terms of \code{square\_of\_four}
as follows:\footnote{Equivalently, we could write

\begin{compactCodeBlock}
>>> flipped_pairs = square_of_four(
...     identity, flip_vert, identity, flip_vert)
\end{compactCodeBlock}
}

\begin{codeBlock}
>>> def flipped_pairs(painter):
...     combine4 = square_of_four(identity, flip_vert,
...                               identity, flip_vert)
...     return combine4(painter)
\end{codeBlock}

\noindent
and \code{square\_limit} can be expressed as\footnote{\code{rotate180} rotates
a painter by 180 degrees (see \link{Exercise 2.50}).  Instead of
\code{rotate180} we could say \code{compose(flip\_vert, flip\_horiz)}, using
the \code{compose} procedure from \link{Exercise 1.42}.}

\begin{codeBlock}
>>> def square_limit(painter, n):
...     combine4 = square_of_four(flip_horiz,
...                               identity,
...                               rotate180,
...                               flip_vert)
...     return combine4(corner_split(painter, n))
\end{codeBlock}

\begin{exerciseGroup}
\phantomsection\label{Exercise 2.45}
\exerciseHeading{Exercise 2.45:} \code{right\_split} and
\code{up\_split} can be expressed as instances of a general splitting
operation.
Define a procedure \code{split} with the property that evaluating

\begin{codeBlock}
>>> right_split = split(beside, below)
>>> up_split = split(below, beside)
\end{codeBlock}

\noindent
produces procedures \code{right\_split} and \code{up\_split} with the same
behaviors as the ones already defined.
\end{exerciseGroup}

\subsubsection*{Frames}

Before we can show how to implement painters and their means of combination, we
must first consider frames.  A frame can be described by three vectors---an
origin vector and two edge vectors.  The origin vector specifies the offset of
the frame's origin from some absolute origin in the plane, and the edge vectors
specify the offsets of the frame's corners from its origin.  If the edges are
perpendicular, the frame will be rectangular.  Otherwise the frame will be a
more general parallelogram.

\link{Figure 2.14} shows a frame and its associated vectors.
In accordance with data abstraction, we need not be specific yet about how
frames are represented, other than to say that there is a constructor
\code{make\_frame}, which takes three vectors and produces a frame, and three
corresponding selectors \code{origin\_frame}, \code{edge1\_frame}, and
\code{edge2\_frame} (see \link{Exercise 2.47}).

% figure 2.15
\begin{imageFigure}{A frame is described by three vectors --- an origin
  and two edges.}
  \includegraphics[width=51mm]{assets/figures/chapter_2/figure_2.15a.pdf}
\end{imageFigure}

We will use coordinates in the unit square \( (0 \le x, y \le 1) \) to specify
images.  With each frame, we associate a \newterm{frame coordinate map}, which
will be used to shift and scale images to fit the frame.  The map transforms
the unit square into the frame by mapping the vector \( \mathbf{v} = (x, y) \)
to the vector sum

$$ {\rm Origin(Frame)} + x \cdot {\rm Edge_1(Frame)} + y \cdot {\rm
Edge_2(Frame)}. $$

\noindent
For example, (0, 0) is mapped to the origin of the frame, (1, 1) to the vertex
diagonally opposite the origin, and (0.5, 0.5) to the center of the frame.  We
can create a frame's coordinate map with the following
procedure:\footnote{\code{frame\_coord\_map} uses the vector operations
described in \link{Exercise 2.46} below, which we assume have been implemented
using some representation for vectors.  Because of data abstraction, it doesn't
matter what this vector representation is, so long as the vector operations
behave correctly.}

\begin{codeBlock}
>>> def frame_coord_map(frame):
...     def mapped(v):
...         return add_vect(
...             origin_frame(frame),
...             add_vect(scale_vect(xcor_vect(v),
...                                 edge1_frame(frame)),
...                      scale_vect(ycor_vect(v),
...                                 edge2_frame(frame))))
...     return mapped
\end{codeBlock}

\noindent
Observe that applying \code{frame\_coord\_map} to a frame returns a procedure
that, given a vector, returns a vector.  If the argument vector is in the unit
square, the result vector will be in the frame.  For example,

\begin{codeBlock}
>>> frame_coord_map(a_frame)(make_vect(0, 0))
\end{codeBlock}

\noindent
returns the same vector as

\begin{codeBlock}
>>> origin_frame(a_frame)
\end{codeBlock}

\begin{exerciseGroup}
\phantomsection\label{Exercise 2.46}
\exerciseHeading{Exercise 2.46:} A two-dimensional vector \( \mathbf{v} \)
running from the origin to a point can be represented as a pair consisting of
an \( x \)-coordinate and a \( y \)-coordinate.  Implement a data abstraction
for vectors by giving a constructor \code{make\_vect} and corresponding
selectors \code{xcor\_vect} and \code{ycor\_vect}.  In terms of your selectors
and constructor, implement procedures \code{add\_vect}, \code{sub\_vect}, and
\code{scale\_vect} that perform the operations vector addition, vector
subtraction, and multiplying a vector by a scalar:

$$
\begin{array}{r@{{}={}}l}
	(x_1, y_1) + (x_2, y_2) 	& (x_1 + x_2, y_1 + y_2), \\
	(x_1, y_1) - (x_2, y_2) 	& (x_1 - x_2, y_1 - y_2), \\
	s \cdot (x, y) 			& (sx, sy).
\end{array}
$$


\phantomsection\label{Exercise 2.47}
\exerciseHeading{Exercise 2.47:} Here are two possible
constructors for frames:

\begin{codeBlock}
>>> def make_frame(origin, edge1, edge2):
...     return (origin, edge1, edge2)

>>> def make_frame(origin, edge1, edge2):
...     return (origin, (edge1, edge2))
\end{codeBlock}

For each constructor supply the appropriate selectors to produce an
implementation for frames.
\end{exerciseGroup}

\subsubsection*{Painters}

A painter is represented as a procedure that, given a frame as argument,
produces a particular image shifted and scaled to fit the frame.  That is to
say, if \code{p} is a painter and \code{f} is a frame, then we produce
\code{p}'s image in \code{f} by calling \code{p} with \code{f} as argument.

\begin{notTaste}
Here we depart from the original, and the departure is worth understanding
because it changes what a painter \emph{is}.

The original's painter draws.  It is called for its effect---a procedure
\code{draw-line} puts a line on the screen---and returns nothing anyone uses.
The original can leave \code{draw-line} undefined because in 1985 what it did
depended on the machine in front of you.

Ours returns the segments it would draw, and something else turns those into a
picture.  A painter is a procedure from a frame to a sequence of segments.  We
do this for a practical reason and a better one.  The practical reason is that
the same painter must produce the figures printed in this book and the ones
drawn in the web edition, and a sequence of segments can become either.  The
better reason is that it is what the rest of this chapter has been doing:
nothing is mutated, nothing is sequenced for effect, and combining two painters
turns out to be joining two sequences---which is why \code{beside}, when we
come to it, is shorter here than in the original.

The procedure that turns a painter into a picture is \code{render}, which we
import rather than write:

\begin{codeBlock}
>>> from pysicp.picture import render, WAVE_SEGMENTS
\end{codeBlock}

\noindent
It stands in the same relation to this section as \code{tail\_recursive} does
to \link{Chapter 1}: something the original's language supplies and Python does
not, provided so that the book's code runs.  \code{WAVE\_SEGMENTS} comes with
it---the original shows the \code{wave} painter but never lists the segments
that make it, so ours are supplied the same way.
\end{notTaste}

Then we can create painters for line drawings, such as the \code{wave} painter
in \link{Figure 2.10}, from sequences of line segments as
follows:\footnote{\code{segments\_to\_painter} uses the representation for line
segments described in \link{Exercise 2.48} below.  The original names this
procedure \code{segments->painter}; an arrow is a perfectly good character in a
Scheme name and not in a Python one.  Note also that the original uses
\code{for-each} here, because its painter draws and the results are discarded,
where ours collects the results and so uses a comprehension.}

\begin{codeBlock}
>>> def segments_to_painter(segments):
...     def painter(frame):
...         to_frame = frame_coord_map(frame)
...         return tuple(
...             make_segment(to_frame(start_segment(s)),
...                          to_frame(end_segment(s)))
...             for s in segments)
...     return painter

>>> wave = segments_to_painter(WAVE_SEGMENTS)
\end{codeBlock}

\noindent
The segments are given using coordinates with respect to the unit square.  For
each segment in the list, the painter transforms the segment endpoints with the
frame coordinate map and draws a line between the transformed points.

Representing painters as procedures erects a powerful abstraction barrier in
the picture language.  We can create and intermix all sorts of primitive
painters, based on a variety of graphics capabilities. The details of their
implementation do not matter.  Any procedure can serve as a painter, provided
that it takes a frame as argument and produces something scaled to fit the
frame.\footnote{For example, the original's \code{rogers} painter was
constructed from a gray-level image.  For each point in a given frame, the
\code{rogers} painter determines the point in the image that is mapped to it
under the frame coordinate map, and shades it accordingly.  By allowing
different types of painters, we are capitalizing on the abstract data idea
discussed in \link{Section 2.1.3}, where we argued that a rational-number
representation could be anything at all that satisfies an appropriate
condition.  Here we're using the fact that a painter can be implemented in any
way at all, so long as it draws something in the designated frame.
\link{Section 2.1.3} also showed how pairs could be implemented as procedures.
Painters are our second example of a procedural representation for data.}

\begin{exerciseGroup}
\phantomsection\label{Exercise 2.48}
\exerciseHeading{Exercise 2.48:} A directed line segment in the
plane can be represented as a pair of vectors---the vector running from the
origin to the start-point of the segment, and the vector running from the
origin to the end-point of the segment.  Use your vector representation from
\link{Exercise 2.46} to define a representation for segments with a constructor
\code{make\_segment} and selectors \code{start\_segment} and
\code{end\_segment}.

\phantomsection\label{Exercise 2.49}
\exerciseHeading{Exercise 2.49:} Use \code{segments\_to\_painter}
to define the following primitive painters:

\begin{enumerate}[a.]

\item
The painter that draws the outline of the designated frame.

\item
The painter that draws an ``X'' by connecting opposite corners of the frame.

\item
The painter that draws a diamond shape by connecting the midpoints of the sides
of the frame.

\item
The \code{wave} painter.  Its segments are \code{WAVE\_SEGMENTS}, imported
above; the exercise is to say what remains to be done with them.

\end{enumerate}

\noindent
Draw each of them.  \code{render} takes a painter and a size in points and
returns an SVG document as a string, which may be written to a file and opened
in any browser:

\begin{codeBlock}
>>> from pathlib import Path
>>> Path("x.svg").write_text(render(x_painter, 200))
\end{codeBlock}
\end{exerciseGroup}

\subsubsection*{Transforming and combining painters}

An operation on painters (such as \code{flip\_vert} or \code{beside}) works by
creating a painter that invokes the original painters with respect to frames
derived from the argument frame.  Thus, for example, \code{flip\_vert} doesn't
have to know how a painter works in order to flip it---it just has to know how
to turn a frame upside down: The flipped painter just uses the original
painter, but in the inverted frame.

Painter operations are based on the procedure \code{transform\_painter}, which
takes as arguments a painter and information on how to transform a frame and
produces a new painter.  The transformed painter, when called on a frame,
transforms the frame and calls the original painter on the transformed frame.
The arguments to \code{transform\_painter} are points (represented as vectors)
that specify the corners of the new frame: When mapped into the frame, the
first point specifies the new frame's origin and the other two specify the ends
of its edge vectors.  Thus, arguments within the unit square specify a frame
contained within the original frame.

\begin{codeBlock}
>>> def transform_painter(painter, origin, corner1, corner2):
...     def transformed(frame):
...         m = frame_coord_map(frame)
...         new_origin = m(origin)
...         return painter(make_frame(
...             new_origin,
...             sub_vect(m(corner1), new_origin),
...             sub_vect(m(corner2), new_origin)))
...     return transformed
\end{codeBlock}

\noindent
Here's how to flip painter images vertically:

\begin{codeBlock}
>>> def flip_vert(painter):
...     return transform_painter(
...         painter,
...         make_vect(0.0, 1.0),    # new origin
...         make_vect(1.0, 1.0),    # new end of edge1
...         make_vect(0.0, 0.0))    # new end of edge2
\end{codeBlock}

\noindent
Using \code{transform\_painter}, we can easily define new transformations.
For example, we can define a painter that shrinks its image to the
upper-right quarter of the frame it is given:

\begin{codeBlock}
>>> def shrink_to_upper_right(painter):
...     return transform_painter(
...         painter, make_vect(0.5, 0.5),
...         make_vect(1.0, 0.5), make_vect(0.5, 1.0))
\end{codeBlock}

\noindent
Other transformations rotate images counterclockwise by 90
degrees\footnote{\code{rotate90} is a pure rotation only for square frames,
because it also stretches and shrinks the image to fit into the rotated frame.}

\begin{codeBlock}
>>> def rotate90(painter):
...     return transform_painter(painter,
...                              make_vect(1.0, 0.0),
...                              make_vect(1.0, 1.0),
...                              make_vect(0.0, 0.0))
\end{codeBlock}

\noindent
or squash images towards the center of the frame:\footnote{The original's
diamond-shaped images were created with \code{squash\_inwards} applied to
\code{wave} and \code{rogers}.}

\begin{codeBlock}
>>> def squash_inwards(painter):
...     return transform_painter(painter,
...                              make_vect(0.0, 0.0),
...                              make_vect(0.65, 0.35),
...                              make_vect(0.35, 0.65))
\end{codeBlock}

\noindent
Frame transformation is also the key to defining means of combining two or more
painters.  The \code{beside} procedure, for example, takes two painters,
transforms them to paint in the left and right halves of an argument frame
respectively, and produces a new, compound painter.  When the compound painter
is given a frame, it calls the first transformed painter to paint in the left
half of the frame and calls the second transformed painter to paint in the
right half of the frame:

\begin{codeBlock}
>>> def beside(painter1, painter2):
...     split_point = make_vect(0.5, 0.0)
...     paint_left = transform_painter(
...         painter1, make_vect(0.0, 0.0),
...         split_point, make_vect(0.0, 1.0))
...     paint_right = transform_painter(
...         painter2, split_point,
...         make_vect(1.0, 0.0), make_vect(0.5, 1.0))
...     def painted(frame):
...         return paint_left(frame) + paint_right(frame)
...     return painted
\end{codeBlock}

\noindent
Observe how the painter data abstraction, and in particular the representation
of painters as procedures, makes \code{beside} easy to implement.  The
\code{beside} procedure need not know anything about the details of the
component painters other than that each painter will produce something in its
designated frame.  Note the last two lines.  Where the original calls the two
painters one after the other for their effects, ours joins what they produce
with \code{+}---the same \code{+} that has been joining tuples since
\link{Section 2.2.1}.  That \code{beside} turns out to be concatenation is the
closure property of \link{Section 2.2} showing up one level higher: painters
compose because the things they produce compose.

\begin{exerciseGroup}
\phantomsection\label{Exercise 2.50}
\exerciseHeading{Exercise 2.50:} Define the transformation
\code{flip\_horiz}, which flips painters horizontally, and transformations that
rotate painters counterclockwise by 180 degrees and 270 degrees.

\phantomsection\label{Exercise 2.51}
\exerciseHeading{Exercise 2.51:} Define the \code{below} operation
for painters.  \code{below} takes two painters as arguments.  The resulting
painter, given a frame, produces the first painter's image in the bottom of the
frame and the second painter's image in the top.  Define \code{below} in two
different ways---first by writing a procedure that is analogous to the
\code{beside} procedure given above, and again in terms of \code{beside} and
suitable rotation operations (from \link{Exercise 2.50}).
\end{exerciseGroup}

\subsubsection*{Levels of language for robust design}

The picture language exercises some of the critical ideas we've introduced
about abstraction with procedures and data.  The fundamental data abstractions,
painters, are implemented using procedural representations, which enables the
language to handle different basic drawing capabilities in a uniform way.  The
means of combination satisfy the closure property, which permits us to easily
build up complex designs.  Finally, all the tools for abstracting procedures
are available to us for abstracting means of combination for painters.

We have also obtained a glimpse of another crucial idea about languages and
program design.  This is the approach of \newterm{stratified design}, the
notion that a complex system should be structured as a sequence of levels that
are described using a sequence of languages.  Each level is constructed by
combining parts that are regarded as primitive at that level, and the parts
constructed at each level are used as primitives at the next level.  The
language used at each level of a stratified design has primitives, means of
combination, and means of abstraction appropriate to that level of detail.

Stratified design pervades the engineering of complex systems.  For example, in
computer engineering, resistors and transistors are combined (and described
using a language of analog circuits) to produce parts such as and-gates and
or-gates, which form the primitives of a language for digital-circuit
design.\footnote{\link{Section 3.3.4} describes one such language.} These parts
are combined to build processors, bus structures, and memory systems, which are
in turn combined to form computers, using languages appropriate to computer
architecture.  Computers are combined to form distributed systems, using
languages appropriate for describing network interconnections, and so on.

As a tiny example of stratification, our picture language uses primitive
elements (primitive painters) that are created using a language that specifies
points and lines to provide the lists of line segments for
\code{segments\_to\_painter}, or the shading details for a painter like
\code{rogers}.  The bulk of our description of the picture language focused on
combining these primitives, using geometric combiners such as \code{beside} and
\code{below}.  We also worked at a higher level, regarding \code{beside} and
\code{below} as primitives to be manipulated in a language whose operations,
such as \code{square\_of\_four}, capture common patterns of combining geometric
combiners.

Stratified design helps make programs \newterm{robust}, that is, it makes it
likely that small changes in a specification will require correspondingly small
changes in the program.  For instance, suppose we wanted to change the image
based on \code{wave} shown in \link{Figure 2.9}.  We could work at the lowest
level to change the detailed appearance of the \code{wave} element; we could
work at the middle level to change the way \code{corner\_split} replicates the
\code{wave}; we could work at the highest level to change how
\code{square\_limit} arranges the four copies of the corner.  In general, each
level of a stratified design provides a different vocabulary for expressing the
characteristics of the system, and a different kind of ability to change it.

\begin{exerciseGroup}
\phantomsection\label{Exercise 2.52}
\exerciseHeading{Exercise 2.52:} Make changes to the square limit
of \code{wave} shown in \link{Figure 2.9} by working at each of the levels
described above.  In particular:

\begin{enumerate}[a.]

\item
Add some segments to the primitive \code{wave} painter of \link{Exercise 2.49}
(to add a smile, for example).

\item
Change the pattern constructed by \code{corner\_split} (for example, by using
only one copy of the \code{up\_split} and \code{right\_split} images instead of
two).

\item
Modify the version of \code{square\_limit} that uses \code{square\_of\_four} so
as to assemble the corners in a different pattern.  (For example, you might
make the large figures face outward from each corner of the square rather than
inward.)

\end{enumerate}
\end{exerciseGroup}

